\documentclass[12pt]{article}
\usepackage[utf8]{inputenc}
\usepackage[T1]{fontenc}
\usepackage{hyperref}
\title{Liberia Kpelle Language Resources for K-12 Curriculum Development}
\author{Compiled by Claude Sonnet 5.5 (claude-sonnet-5-5) for Andr\'as Kornai}
\date{}
\begin{document}
\maketitle

\section{Introduction}
Liberia Kpelle (ISO 639-3 xpe, Glottocode libe1247) is a Mande language of Liberia with roughly one million speakers. Sources found in search disagree on whether it is taught in schools. One university guide says it is not officially taught in Liberian schools and that Liberia has no official mother-tongue instruction policy (\url{https://wisc.pb.unizin.org/lctlresources/chapter/about-kpelle/}). Wikipedia says it is taught (\url{https://en.wikipedia.org/wiki/Kpelle_language}). This report treats the question as unresolved. Published materials are few: one small dictionary, a few older reader and grammar items, a complete Bible, and an indigenous syllabary. No ready-made K-12 curriculum turned up in the searches.

\section{Primary Educational Resources}
\begin{itemize}
\item Kpelle reader \emph{Gbɛlɛɛ-kɔlɔi 2b} (1960), a translation of a Bible primer from the Kpelle Literacy Center, Totota. It is catalogued in the Indiana University African Ephemera collection: \url{https://collections.libraries.indiana.edu/africancollections/items/show/7405}.
\item Better Jallah's rural church literacy classes, which use Kpelle Bibles, are described by Lutheran Bible Translators: \url{https://lbt.org/better-jallahs-vision-for-liberia/}. They show that adult literacy work in Kpelle exists, though rural access is limited.
\end{itemize}

\section{Linguistic and Reference Materials}
\begin{itemize}
\item Elizabeth Winkler's 1997 Kpelle-English dictionary (Bong County Kpelle) has about 1,100 entries. A review is at \url{https://linguistlist.org/issues/10/68/}.
\item Winkler's online Kpelle dictionary project at Western Kentucky University is described in a sabbatical report, \url{https://digitalcommons.wku.edu/sabb_rpt/4}, and in a news article, \url{https://wkuherald.com/48111/life/professor-helps-document-african-language-for-the-first-time}.
\item The 2005 Kpelle-English Dictionary (Leidenfrost and McKay, Lutheran Church in Liberia) and other items are listed in the WEBBALL bibliography: \url{https://reflex.cnrs.fr/Lexiques/webball/biblio_bas.php?id=24606}.
\item Westermann and Melzian (1930) is recorded at \url{https://wals.info/refdb/record/Westermann-and-Melzian-1930}.
\item The Kpelle syllabary, reportedly devised around 1935, is described at \url{https://en.wikipedia.org/wiki/Kpelle_syllabary}.
\item The overview and teaching-resource guide for less commonly taught languages is at \url{https://wisc.pb.unizin.org/lctlresources/chapter/about-kpelle/}.
\end{itemize}

\section{Religious and Community Resources}
\begin{itemize}
\item A complete Kpelle Bible was dedicated in 2014, according to the Lutheran Bible Translators press release at \url{https://lbt.org/?p=12077}. Many rural churches still lack copies.
\item The Joshua Project lists the Kpelle of Liberia as having a complete Bible: \url{https://www.joshuaproject.net/pray/cards/ROG3/LI/Ctry/PeopNameInCountry}.
\item The Lutheran Church in Liberia's earlier translation effort is described at \url{https://blogs.elca.org/handinhand/?p=2864}.
\item A HathiTrust record exists for a Kpelle Genesis: \url{https://catalog.hathitrust.org/Record/011809845}.
\end{itemize}

\section{Digital and Audio Resources}
\begin{itemize}
\item Global Recordings Network lists Kpelle audio: \url{https://www.globalrecordings.net/en/language/kpe}.
\item The LBT release mentions recorded Scripture on Proclaimer devices and SD cards (\url{https://lbt.org/?p=12077}).
\end{itemize}

\section{Curriculum Development Context}
Native-language materials for school use are scarce. The foundations are the Bible, the small dictionaries, the older reader, the syllabary and church literacy programmes. A curriculum would first need a standard orthography decision (Latin or syllabary), graded readers, and a larger lexicon. Kpelle also has speakers in Guinea, and Guinean Kpelle materials were not searched here. The Lutheran Church in Liberia and its translation centre are the obvious partners.

\end{document}